\Page{13}{\textenglish{03 / CAMPAIGNS}}{התאומות של \textenglish{H\&M}}{45}{52}

\FeatureIntro{PaperPink}{בתחילת יולי 2025 פרסמה \textenglish{H\&M} סדרת דנים שבה מופיעות תאומות דיגיטליות של דוגמניות אמיתיות, על רקעים שמזכירים בירות אופנה. כשמותג רוצה שהקהל יזהה פנים וטון משותף בכמה סצנות, תאומה דיגיטלית נותנת בסיס להעמיד את אותה דמות בהקשרים חזותיים שונים. ההשקה כללה סרטון מאחורי הקלעים. בהודעת המותג הופיעו גם הצלם ג׳וני קנגסניימי והדוגמנית ונסה מודי, חלק מן הצוות שעבד על המהלך. \SourceRef{8}}

\CaseLabel{H\&M}{https://hmgroup.com/news/hm-continues-its-exploration-of-creativity-with-ai/}{לסדרה ולסרטון ההפקה}

{\fontsize{14.5}{21.7}\selectfont

\StoryColumns{את ההכנה אפשר לראות בדיווח של \textenglish{The Business of Fashion} מחודש מרץ. לפי הכתבה, כדי ליצור תאומה, תיעדו את הדוגמנית בצילום מזוויות שונות, בתאורות שונות ובתנועה. החומר נתן לצוות בסיס לבנות ייצוג שאפשר לחזור אליו. כמה חודשים לאחר מכן הגיעו התמונות הראשונות. זו תוצאה שונה מיצירת דמות שאין לה מקבילה אנושית. חלק מהעניין עבור הצוות הוא לשמור על מראה מזוהה כשהתמונה נבנית מחדש. \SourceRef{8}}{גם חלוקת העבודה משתנה. דוגמנית יכולה להצטלם כדי להכין את הייצוג, וצלם יכול להיות מעורב גם כשהפריים האחרון לא נוצר בצילום ישיר. בצד זה נמצאים הקריאייטיב, מי שמכינים גרסאות והגורמים שמאשרים כיצד ייעשה השימוש הבא בדמות. התיעוד הפומבי מאפשר לעקוב אחר המעבר מחומר מצולם לפרסום. פרטי החוזים והתגמול של כל אחת מן המשתתפות אינם זמינים בו. \SourceRef{8}}

\par}

\vspace{5mm}

{\fontsize{14.5}{21.9}\selectfont

ארגון \textenglish{Model Alliance} הגיב עוד במרץ והעלה חששות לגבי ההגנות על דוגמניות וצוותים שהעבודה שלהם עשויה להשתנות. המקרה מסביר למה הדיון קדם לתמונות. תאומה היא נכס שאפשר להפיק ממנו גרסה חדשה גם אחרי יום הצילום הראשון. בשביל ההפקה זו גמישות. לדוגמנית חשוב לדעת אילו שימושים אישרה, ולמי יש סמכות לעצור גרסה שחורגת מן ההסכמה המקורית. \SourceRef{8}

\par}
